\BourbakiExerciseGroup

¶ 1) Let R be an open covering of a compact space X. Show that there is an entourage V of the uniformity of X such that, for each \(x \in X\) , \(V(x)\) is contained in a set belonging to R {[}remark that for each \(x \in X\) there is an entourage \(W_x\) such that \(\dot{W}_x(x)\) is contained in a set of R, and cover X by a finite number of sets \(W_x(x)\) {]}.

\begin{enumerate}
\def\labelenumi{\arabic{enumi})}
\setcounter{enumi}{1}
\item
  Prove that a quasi-compact space X is uniformizable if and only if its topology is the inverse image of the topology of a compact space Y under a surjective mapping \(f: X \rightarrow Y\) .
\item
  Let \(f\) be a continuous mapping of a compact space \(\mathbf{X}\) into a compact space \(\mathbf{Y}\) , and let \(\mathbf{V}\) be an open entourage of \(\mathbf{X}\) such that the relation \(f(x) = f(y)\) implies \((x, y) \in \mathbf{V}\) . Show that there is an entourage \(W\) of \(\mathbf{Y}\) such that the relation \((f(x), f(y)) \in W\) implies \((x, y) \in V\) .
\item
  Let X be the locally compact space defined in Chapter I, § 9, Exercise II. Show that every neighbourhood of the diagonal \(\Delta\) in \(X \times X\) contains a set of the form \([x, b[ \times [x, b[\) . Deduce that there is only one uniformity on X compatible with the topology of X, and that the completion of X for this uniformity can be identified with the interval \(X' = [a, b]\) carrying the topology \(\mathcal{C}v(X)\) . (Observe that every ultrafilter on X must be a Cauchy filter with respect to any uniformity compatible with the topology of X).
\item
  Show that a uniform space X is precompact if and only if every ultrafilter on X is a Cauchy filter. (Remark that if X is not precompact then there is an entourage V of X such that the sets \(\mathbb{C}\mathrm{V}(x)\) , as x runs through X, generate a filter on X).
\item
  Let X be a uniform space in which every sequence of points has at least one cluster point. Show that X is precompact (use reductio ad absurdum).
\item
  A subset \(\mathbf{A}\) of a uniform space \(\mathbf{X}\) is said to be bounded if for each entourage \(\mathbf{V}\) of \(\mathbf{X}\) there is a finite set \(\mathbf{F}\) and an integer \(n > 0\) such that \(\mathbf{A} \subset \overset{n}{\mathbf{V}}(\mathbf{F})\) .
\end{enumerate}

\begin{enumerate}
\def\labelenumi{\alph{enumi})}
\item
  The union of two bounded subsets is bounded. The closure of a bounded subset is bounded. Every precompact subset is bounded.
\item
  If \(f: \mathbf{X} \to \mathbf{Y}\) is uniformly continuous, then the image under \(f\) of every bounded subset of \(\mathbf{X}\) is a bounded subset of \(\mathbf{Y}\) .
\item
  In a product of non-empty uniform spaces, a set is bounded if and only if each of its projections is bounded.
\item
  For each symmetric entourage V of X, let \(R_{V}\) denote the union of the sets \(\stackrel{n}{V}\) as n runs through the integers \textgreater0. Show that \(R_{V}\) is both open and closed in \(X \times X\) and that, in the notation of no. 4,
\end{enumerate}

\[
\mathbf {R} _ {\mathbf {v}} (x) = \mathbf {A} _ {x, \mathbf {v}}
\]

for each \(x \in \mathbf{X}\) . Suppose that \(\mathbf{R}_{\mathbf{V}} = \mathbf{X} \times \mathbf{X}\) for each symmetric entourage \(\mathbf{V}\) of \(\mathbf{X}\) (this will be the case in particular if \(\mathbf{X}\) is connected); show that a set \(\mathbf{A} \subset \mathbf{X}\) is bounded if and only if, for each symmetric entourage \(\mathbf{V}\) of \(\mathbf{X}\) and each \(x_0 \in \mathbf{X}\) , there is an integer \(n > 0\) such that \(\mathbf{A} \subset \mathbf{V}(x_0)\) .

\begin{enumerate}
\def\labelenumi{\arabic{enumi})}
\setcounter{enumi}{7}
\item
  Let X be a uniform space, V a closed entourage of X, A a compact subset of X. Show that V(A) is closed in X (cf.~Chapter I, § 10, no. 2, Theorem 1, Corollary 5).
\item
  Let X be a locally compact space which carries a uniformity compatible with its topology and is such that there is an entourage V with the property that \(V(x)\) is relatively compact in X for all \(x \in X\) . a) Show that X is complete with respect to this uniformity.
\end{enumerate}

\begin{enumerate}
\def\labelenumi{\alph{enumi})}
\setcounter{enumi}{1}
\item
  Let \(\mathbf{U}\) be a symmetric entourage of \(\mathbf{X}\) such that \(\dot{\mathbf{U}}\subset\mathbf{V}\) . Show that if \(\mathbf{A}\) is relatively compact in \(\mathbf{X}\) , then so is \(\mathbf{U}(\mathbf{A})\) .
\item
  Show that X is paracompact {[}use b), Exercise 7 and Chapter I, § 9, no. 10, Theorem 5{]}.
\item
  Let W be a closed symmetric entourage of X such that \(\dot{W} \subset V\) . Show that if A is closed in X then so is \(\mathbf{W}(\mathbf{A})\) (use Exercise 8).
\end{enumerate}

\begin{enumerate}
\def\labelenumi{\arabic{enumi})}
\setcounter{enumi}{9}
\tightlist
\item
  Let X be a Hausdorff uniform space which has an entourage V such that \(\mathbf{V}(x)\) is precompact for each \(x \in X\) . Show that the completion \(\hat{X}\) is locally compact and satisfies the conditions of Exercise 9. {[}If W is a symmetric open entourage of X such that \(\hat{W} \subset V\) , show that the closure \(\overline{W}\) of W in \(\hat{X} \times \hat{X}\) is such that \(\overline{W}(x)\) is compact for each \(x \in \hat{X}\) {]}.
\end{enumerate}

¶ 11) Let X be a Hausdorff uniform space, U its uniformity, F(X) the set of all non-empty closed subsets of X. Let F(X) carry the uniformity induced by the uniformity U defined in § 1, Exercise 5 a).

\begin{enumerate}
\def\labelenumi{\alph{enumi})}
\item
  Show that if \(X'\) is precompact then so is \(\mathfrak{F}(\mathbf{X})\) . {[}If \((\mathbf{A}_{i})\) is a finite covering of X by V-small sets (V symmetric), show that the sets \(\tilde{\mathbf{V}}(\mathbf{B}_{k})\) cover \(\mathfrak{F}(\mathbf{X})\) , where the \(B_{k}\) are all the unions of sets \(A_{i}]\) .
\item
  A point \(\mathbf{A} \in \mathfrak{F}(\mathbf{X})\) is such that every neighbourhood of \(\mathbf{A}\) in the topology induced by \(\mathfrak{G}(\tilde{\mathcal{U}})\) on \(\mathfrak{F}(\mathbf{X})\) contains a neighbourhood of \(\mathbf{A}\) in the topology induced by \(\mathfrak{G}_{\Theta}\) (Chapter I, § 8, Exercise 12) on \(\mathfrak{F}(\mathbf{X})\) , if and only if \(\mathbf{A}\) is precompact.
\item
  Show that the topologies \(\mathfrak{G}(\mathfrak{U})\) and \(\mathfrak{G}_{\Theta}\) induce the same topology on the set \(\mathfrak{K}(\mathbf{X})\) of compact subsets of \(\mathbf{X}\) {[}use \(b\) ) and Exercise 5 of §1{]}.
\end{enumerate}

\begin{enumerate}
\def\labelenumi{\arabic{enumi})}
\setcounter{enumi}{11}
\tightlist
\item
  \begin{enumerate}
  \def\labelenumii{\alph{enumii})}
  \tightlist
  \item
    Let R be a Boolean algebra, and identify R with a Boolean algebra of subsets of the set \(\Omega\) of maximal prefilters on R (Chapter I, § 6, Exercise 20). For each finite partition \(\overline{\omega} = (\mathbf{A}_i)\) of \(\Omega\) , formed of sets belonging to R, let \(\mathbf{C}_{\overline{\omega}}\) denote the subset \(\bigcup_{i} (\mathbf{A}_i \times \mathbf{A}_i)\) of \(\Omega \times \Omega\) .
  \end{enumerate}
\end{enumerate}

Show that the sets \(C_{\overline{\omega}}\) form a fundamental system of entourages of a Hausdorff uniformity on \(\Omega\) . The completion, \(\hat{\Omega}\) , of \(\Omega\) with respect to this uniformity is a totally disconnected compact space. A subset of \(\hat{\Omega}\) is both open and closed in \(\hat{\Omega}\) if and only if it is of the form \(\overline{\mathbf{A}}\) , where \(\mathbf{A} \in \mathbf{R}\) . Show that \(\mathbf{A} \to \overline{\mathbf{A}}\) is a bijective mapping of \(\mathbf{R}\) onto the set of subsets of \(\hat{\Omega}\) which are both open and closed, and that \(\overline{\mathbf{C}}\mathbf{A} = \overline{\mathbf{C}}\overline{\mathbf{A}}\) , \(\overline{\mathbf{A} \cup \mathbf{B}} = \overline{\mathbf{A}} \cup \overline{\mathbf{B}}\) and \(\overline{\mathbf{A} \cap \mathbf{B}} = \overline{\mathbf{A}} \cap \overline{\mathbf{B}}\) . Deduce that the canonical mapping \(\Omega \to \hat{\Omega}\) is bijective.

\begin{enumerate}
\def\labelenumi{\alph{enumi})}
\setcounter{enumi}{1}
\tightlist
\item
  Let X be a Hausdorff topological space in which the sets that are both open and closed form a base for the topology. Take R to be the Boolean algebra formed by subsets of X which are both open and closed; then the topology induced on X by the topology of the completion \(\hat{\mathbf{X}}\) of X with respect to the uniformity \(\mathcal{U}\) defined in \(a\) ) is the given topology on X. Further, the maximal elements of the set of filters on X having a base formed of open sets, and the maximal elements of the set of filters on X having a base formed of closed sets are Cauchy filters with respect to the uniformity \(\mathcal{U}\) . Deduce that there are continuous surjections \(\varphi: \hat{\mathbf{X}} \to \hat{\mathbf{X}}, \psi: \mathbf{X}' \to \hat{\mathbf{X}}\) , where \(\hat{\mathbf{X}}, \mathbf{X}'\) are the spaces defined in Chapter I, § 9, Exercises 25 and 24. * Show that if X is the rational line Q, then \(\psi\) is not bijective.* If X is extremally disconnected (Chapter I, § 11, Exercise 21), then \(\psi\) is bijective, \(\hat{\mathbf{X}}\) is extremally disconnected and can be identified with the semi-regular space associated with \(\mathbf{X}'\) (Chapter I, § 8, Exercise 20); thus extremally disconnected spaces can be characterized as dense subspaces of extremally disconnected compact spaces. Hence give an example of an extremally disconnected compact space which has no isolated points {[}(cf.~Chapter I, § 11, Exercise 21 f){]}.
\end{enumerate}

In particular, if we take X to be a discrete space, then \(\hat{X}\) can be identified with the ultrafilter space of X (Chapter I, § 9, Exercise 26).

13) Let X be a locally compact space and U an open subset of X. Let B be the union of U and the relatively compact components of CU. Show that B is open and is the union of U and those subsets of CU which are both open and compact in CU. (Embed X in its one-point compactification \(X' = X \cup \{\omega\}\) , observe that in \(X' - U\) the component of \(\omega\) is \(\{\omega\} \cup CB\) , and use Proposition 6 of no. 4).

\begin{enumerate}
\def\labelenumi{\arabic{enumi})}
\setcounter{enumi}{13}
\item
  Let X be a compact space, and let B be a filter base on X consisting of connected closed sets. Show that the (non-empty) intersection of the sets of B is closed and connected (argue by reductio ad absurdum, using Proposition 4 of no. 3 and Chapter I, § 9, no. 1, Theorem 1).
\item
  Let X be a compact space. Then the space \(\mathfrak{F}(\mathbf{X})\) of non-empty closed subsets of X, with the topology induced by \(\mathcal{C}(\tilde{\mathcal{U}})\) (Exercise 11), is compact (Chapter I, § 9, Exercise 13). Consider an ultrafilter \(\Psi\) on \(\mathfrak{F}(\mathbf{X})\) , and suppose that for each entourage V of the uniformity of X, and each \(\mathfrak{X} \in \Psi\) , there exists \(\mathbf{M} \in \mathfrak{X}\) such that any two points of \(\mathbf{M}\) can be joined by a V-chain contained in \(\mathbf{M}\) . Show that the limit point A of \(\Psi\) in \(\mathfrak{F}(\mathbf{X})\) is a connected subset of \(\mathbf{X}\) (use Proposition 6 of no. 4).
\item
  Let X be a connected compact space and let A, B be two disjoint non-empty closed subsets of X. Show that there is a component of \(\mathbb{C}(\mathrm{A} \cup \mathrm{B})\) whose closure meets both A and B. (Let U be an entourage of the uniformity of X such that \(\overline{\mathbf{U}}(\mathbf{A}) \cap \overline{\mathbf{U}}(\mathbf{B}) = \varnothing\) . Show first that the set \(\mathfrak{M}_{\mathrm{U}}\) of components of \(\overline{\mathbb{C}(\mathrm{U}(\mathbf{A}) \cup \mathrm{U}(\mathbf{B}))}\) which meet both \(\overline{\mathrm{U}(\mathbf{A})}\) and \(\overline{\mathrm{U}(\overline{\mathrm{B}})}\) is not empty, by proving that for each symmetric entourage W ⊂ U there exist \(x \in \mathrm{U}(\mathbf{A})\) and \(y \in \mathrm{U}(\mathbf{B})\) and a W-chain which joins x to y and all of whose points other than x and y belong to \(\mathbb{C}(\mathrm{U}(\mathbf{A}) \cup \mathrm{U}(\mathbf{B}))\) , and then applying Exercise 15. For each entourage V ⊂ U, show that every K ∈ \(\mathfrak{M}_{\mathrm{V}}\) contains a set H ∈ \(\mathfrak{M}_{\mathrm{U}}\) , and that the set \(\mathfrak{M}_{\mathrm{U},\mathrm{V}}\) of sets H ∈ \(\mathfrak{M}_{\mathrm{U}}\) which are contained in a set K ∈ \(\mathfrak{M}_{\mathrm{V}}\) is closed in \(\mathfrak{F}(\overline{\mathrm{X}})\) ; and conclude that as V runs through the set of entourages contained in U, the intersection of the sets \(\mathfrak{M}_{\mathrm{U},\mathrm{V}}\) is not empty).
\item
  Let X be a connected locally compact space. Let K be any non-empty compact subset of X distinct from X. Show that every component of K meets the frontier of K in X (use the Corollary to Proposition 6 of no. 4). Deduce that if A is any non-empty relatively compact open subset of X distinct from X, then the closure of every component of A meets CA.
\end{enumerate}

¶ 18) a) Show that a compact connected space X cannot be the union of a countable infinity of mutually disjoint non-empty closed sets {[}by reductio ad absurdum: if (F \(_{n}\) ) is a countably infinite partition of X into closed sets, show with the help of Exercise 17 that there is a compact connected set K which does not meet F \(_{1}\) but meets an infinite number of sets F \(_{n}\) ; in order to do this consider the component of a point of F \(_{2}\) with respect to a compact neighbourhood of F \(_{2}\) not meeting F \(_{1}\) . Then use induction{]}.

\begin{enumerate}
\def\labelenumi{\alph{enumi})}
\setcounter{enumi}{1}
\tightlist
\item
  Extend the result of a) to a locally compact connected space X under one or the other of the following two supplementary conditions: one of the \(F_{n}\) is compact and connected, or X is locally connected. {[}In the first case, reduce to a) by means of Exercise 17{]}.
\end{enumerate}

\(^{*}c)\) Let \(\mathbf{X}\) be the subspace of \(\mathbb{R}^3\) which is the union of the following subspaces: \(A_n\) is the half-line \(x > 0\) , \(y = 1/n\) , \(z = 0\) ( \(n \geqslant 1\) ); \(B_n\) is the interval \(2n < x < 2n + 2\) , \(y = 0\) , \(z = 0\) ( \(n \geqslant 0\) ); \(C_n\) is the set defined by the relations

\[
x = 2 n + 1, \quad 0 \leqslant y \leqslant \frac {1}{n + 1}, \quad z = y \left(y - \frac {1}{n + 1}\right),
\]

where \(n \geqslant 0\) .

Show that X is connected and locally compact and is the union of a countable infinite family of mutually disjoint connected non-empty closed sets. *

¶ 19) A compact connected space X is said to be irreducible between two of its points x, y if there is no compact connected subset of X other than X itself which contains both x and y.

\begin{enumerate}
\def\labelenumi{\alph{enumi})}
\item
  If \(x\) and \(y\) are two distinct points of a compact connected space \(X\) , show that \(X\) has a compact connected subspace \(K\) which is irreducible between \(x\) and \(y\) (use Exercise 14 and Zorn's Lemma).
\item
  Show that a compact connected space having at least two points cannot be irreducible between every pair of its points (use Exercise 17).
\item
  Let X be a compact connected space, and suppose that there is a point \(a \in X\) which has a connected closed neighbourhood \(V \neq X\) . If x and y are any two points of CV, show that there is a compact connected subset of X which contains x and one of the points a, y but not the other (use Exercise 17). Under the same hypotheses, show that if x, y, z are any three distinct points of X, then X cannot be simultaneously irreducible between x and y, between y and z, and between x and z.
\end{enumerate}

\begin{enumerate}
\def\labelenumi{\arabic{enumi})}
\setcounter{enumi}{19}
\tightlist
\item
  Let X be a compact connected space and let L be the set of all points of X which have a fundamental system of connected neighbourhoods. The points of S = CL are called the singular points of X. The points of the interior \(\dot{L}\) of L and the components of \(\bar{S}\) are called the prime constituents of X.
\end{enumerate}

\begin{enumerate}
\def\labelenumi{\alph{enumi})}
\item
  Let \(\mathbf{A} \neq \mathbf{X}\) be a non-empty open subset of \(\mathbf{X}\) , let \(\mathbf{K}\) be a component of \(\overline{\mathbf{A}}\) and let \(\mathbf{F}\) be the closure of \(\overline{\mathbf{A}} \cap \mathbb{C}\mathbf{K}\) . Show that if \(\mathbf{A} \cap \mathbf{K} \cap \mathbf{F}\) is not empty then all its points are singular. If \(x \in \mathbf{A} \cap \mathbf{F} \cap \mathbf{K}\) , the component \(\mathbf{Q}\) of \(x\) in \(\mathbf{F}\) meets \(\mathbb{C}\mathbf{A}\) {[}consider a point of \(\mathbb{C}\mathbf{K}\) contained in \(V(x)\) and its component in \(\overline{\mathbf{A}}\) , and show by using Exercise 17 that the set of points of \(\mathbf{F}\) which can be joined to \(x\) by a V-chain in \(\mathbf{F}\) meets \(\mathbb{C}\mathbf{A}]\) . Deduce that the prime constituent \(\mathbf{P}\) which contains \(x\) meets \(\mathbb{C}\mathbf{A}\) (remark that \(\mathbf{P}\) contains the component of \(x\) in \(Q \cap A\) , and use Exercise 17 applied to the compact connected space \(Q\) and the open subset \(Q \cap A\) of \(Q\) ).
\item
  Deduce from a) that if P is a prime constituent of X and U is a neighbourhood of P, then P has a connected neighbourhood contained in U (by reductio ad absurdum).
\item
  Let V be a symmetric entourage of X, and let \(V'\) denote the set of pairs \((x, y)\) of points of X which can be joined by a V-chain, all of whose points are singular, except perhaps for x and y. As V runs through the symmetric entourages of X, the corresponding sets \(V'\) form a fundamental system of entourages of a uniformity (in general not Hausdorff) on X. Let \(X'\) be the associated Hausdorff space, and show that the inverse images of the points of \(X'\) under the canonical mapping \(X \rightarrow X'\) are the prime constituents of X. The (compact connected) space \(X'\) is called the space of prime constituents of X. Show that \(X'\) is locally connected {[}use b) and Chapter I, § 10, no. 4, Proposition 8{]}.
\end{enumerate}

\(^{*}d)\) Let \((r_{n})\) be the sequence of all the rational numbers contained in \(]0,1[\) , arranged in some order. For each irrational \(x\in[0,1]\) , let

\[
f (x) = \sum_ {n} 2 ^ {- n} \sin 1 / (x - r _ {n}).
\]

Let X be the closure in \(R^{2}\) of the graph of f.~Show that X is connected and irreducible between its points with abscissas 0 and 1, but that X has only one prime constituent.*

*21) In the locally compact space X defined in Exercise 20 b) of Chapter I, § 10, let S be the equivalence relation whose classes are the components of X. Show that the quotient space X/S is not Hausdorff.*

\begin{enumerate}
\def\labelenumi{\arabic{enumi})}
\setcounter{enumi}{21}
\tightlist
\item
  Show that every totally disconnected compact space X is homeomorphic to the inverse limit of an inverse system of finite discrete spaces. (Consider the finite partitions of X into open sets and use the Corollary to Proposition 6 of no. 4).
\end{enumerate}

¶ 23) Let \(f: \mathbf{X} \to \mathbf{Y}\) be a continuous open mapping, \(\mathbf{X}\) locally compact and \(\mathbf{Y}\) Hausdorff.

\begin{enumerate}
\def\labelenumi{\alph{enumi})}
\item
  If K is any compact connected subset of Y and C is any compact component of \(\overline{f}^{1}(K)\) , show that \(f(C)=K\) (reduce to the case K=Y and use the Corollary to Proposition 6 of no. 4). *Give an example of a non-compact component \(C'\) of \(\overline{f}^{1}(K)\) such that \(f(C')\neq K_*\)
\item
  Suppose in addition that Y is locally compact and locally connected. Show that if U is any connected open subset of Y and R is any relatively compact component of \(\overline{f}^{1}(U)\) in X, then \(f(R)=U\) . {[}Observe that in a locally compact locally connected space, a connected open set U is the union of the compact connected sets contained in U and containing a given point \(y_{0}\in U\) ; then apply a).{]}
\item
  Suppose in addition that X is locally connected. Show that if K is any connected subset of Y which is either open or compact with a non-empty interior, and if H is any compact subset of X, then there is only a finite number of components of \(\overline{f}^{1}(K)\) contained in H. {[}Use a) and b){]}.
\end{enumerate}

